% - GitHub
%     - Co je GitHub (hostovaná Git služba, největší platforma pro open-source vývoj)
%       - Stručná historie, jak vznikl, kdy byl odkoupen Microsoftem
%       - Jak se liší od prostého gitu (sociální funkce, webové rozhraní, CI/CD, Issues, wiki...)

\section{Úvod do GitHubu}
\label{sec:github_introduction}
GitHub je jedna z~mnoha hostovaných služeb pro správu Git repozitářů. Rozšiřuje základní funkcionalitu Gitu o~sociální funkce, webové rozhraní, \CICD a~další.
Nejznámější alternativy jsou \devTool{GitLab}{https://gitlab.com}, který je na rozdíl od GitHubu open-source a~umožňuje vlastní provoz (self-hosting), a~\devTool{Bitbucket}{https://bitbucket.org}~\cite{psl_git_tutorial_main}.

Po vzniku Gitu pro jádro operačního systému Linux se brzy ukázalo, že vlastní hostování repozitářů není vždy dostačující a~vznikla poptávka po sdílené cloudové službě. V~roce 2008 proto vznikl GitHub a~od té doby se stal největší platformou pro open-source vývoj. V~roce 2018 byl odkoupen společností Microsoft~\cite{psl_git_tutorial_main}.

Hlavní rozdíl mezi Gitem a~GitHubem spočívá v~tom, že Git je nástroj pro správu verzí (\VCS), zatímco GitHub je kompletní vývojová platforma. Kromě samotného Gitu nabízí možnosti \CICD a~nasazení, sledování chyb a~požadavků (issue tracking), správu pull requestů, diskuzní fóra komunity a~další sociální funkce. GitHub také umožňuje hostování statických stránek prostřednictvím služby \devTool{GitHub Pages}{https://pages.github.com}, která se využívá například pro hostování dokumentace projektů. Samozřejmostí je webové rozhraní pro veškerou správu včetně úpravy kódu a~kontroly pull requestů.
